% ===================================================================
%  図表第 12 号 — 停車場。— 鉄道。— 船。
%
%  Ceci n'est PAS une transcription : c'est une traduction, faite en
%  2026, en japonais d'aujourd'hui. Voir texto/ja/10-zuhyo-01.tex
%  pour la consigne, qui vaut pour les seize fichiers.
% ===================================================================

%%K t12-apar-1 apar
\VUsaut{4.0mm}
\VUtitre{15.0pt}{60}{図表第 12 号}
\VUsaut{2.4mm}
\VUpk{11.4pt}{40}{停車場。--- 鉄道。--- 船。}
\VUsaut{3.4mm}
\textsc{注。} --- \textit{国ごとに用いる時刻が違うので、例として}
\VUgras{フランスの掛図}\textit{の時刻を取った。}

%%K t12-01-1 p
1. --- 私はいまドイツを一巡りしてきた。発つ前に汽車の出る時刻を知ろうと
\VUgras{時刻表} \textsuperscript{(15)} を一冊買った。一番便のよい汽車が夜九時に
パリを出て一時十三分にストラスブールに着くと分かったので、身支度をし、翌日
車を呼んで停車場まで運ばせた。

%%K t12-02-1 p
2. --- 田舎の老いた\VUgras{司祭} \textsuperscript{(2)} の後ろから大きな出発の間に
入ったとき --- 司祭は手に\VUgras{旅の鞄} \textsuperscript{(3)} を提げていた ---
もう多くの\VUgras{旅客} \textsuperscript{(1)} が来ていた。

%%K t12-02-2 p
私は\VUgras{電報} \textsuperscript{(5)} を打たねばならなかったので、\VUgras{検査係}
\textsuperscript{(6)} に\VUgras{電信の室} \textsuperscript{(4)} を尋ねた。

%%K t12-03-1 p
3. --- \VUgras{待合室} \textsuperscript{(7)} へ行く前に、往復の\VUgras{切符}
\textsuperscript{(9)} を買いに行った。

%%K t12-04-1 p
4. --- \VUgras{切符売場} \textsuperscript{(8)} ではあまり待たずにすんだ。人が
少なかったからである。休みで発つ\VUgras{兵} \textsuperscript{(10)} が親切に順を
譲ってくれたので、\VUgras{駅長} \textsuperscript{(13)} に幾つか尋ねる暇もできた。
駅長は\VUgras{警察の署長} \textsuperscript{(12)} と当番の\VUgras{憲兵}
\textsuperscript{(11)} と話していた。

%%K t12-tit-1 sub
\VUpk{10.2pt}{40}{荷を託すこと。}
\VUsaut{3.0mm}

%%K t12-05-1 p
5. --- \VUgras{本と新聞の売店} \textsuperscript{(14)} で新聞を一部買ってから、
自分の\VUgras{荷} \textsuperscript{(17)} を量らせた。\VUgras{赤帽}
\textsuperscript{(16)} がいま\VUgras{荷車} \textsuperscript{(19)} でそれを運んで
きたのである。\VUgras{係の人} \textsuperscript{(22)} が\VUgras{秤} \textsuperscript{(21)}
を受けもっていて、私の行李は三十五キロ、五キロの\VUgras{超過}になるから
\VUgras{荷の窓口} \textsuperscript{(23)} で割増を払わねばならないと言った。

%%K t12-06-1 p
6. --- 早く来ていたので、停車場の旅客の行き来を見る暇があった。ある者は
\VUgras{預り所} \textsuperscript{(25)} に小さな包を預けたり受け取ったりし、ある者は
\VUgras{時刻表} \textsuperscript{(26)} を見ていた。着いた旅客は\VUgras{出口}
\textsuperscript{(32)} へ急ぎ、手に自分の\VUgras{手提} \textsuperscript{(24)} を提げて
いた。何人かは\VUgras{荷の渡し口} \textsuperscript{(27)} で待ち、自分の
\VUgras{荷札} \textsuperscript{(29)} を出して行李や\VUgras{木箱} \textsuperscript{(28)}
や\VUgras{籠} \textsuperscript{(30)} を受け取っていた。

%%K t12-07-1 p
7. --- \VUgras{収税の役人} \textsuperscript{(33)} が包を念入りに調べ、申告すべき
ものがないかを見ていた。

%%K t12-08-1 p
8. --- ある旅客は\VUgras{食堂} \textsuperscript{(34)} へ行き、\VUgras{給仕}
\textsuperscript{(35)} が忙しく世話をしていた。急がねばならない。その\VUgras{汽車}
\textsuperscript{(36)} は急行で、停車場に長くは止まらないからである。

%%K t12-09-1 p
9. --- それから出発の\VUgras{ホーム}へ行き、乗り換えずにすむ直通の
\VUgras{車輌} \textsuperscript{(37)} を探した。廊下のついた車に乗った。中には
煙草を吸う\VUgras{仕切} \textsuperscript{(38)} があり、もう一つは「婦人だけ」に
取ってある。

%%K t12-tit-2 sub
\VUpk{10.2pt}{40}{夜の出発。}
\VUsaut{3.0mm}

%%K t12-10-1 p
10. --- \VUgras{踏段} \textsuperscript{(40)} を上がり、\VUgras{扉}
\textsuperscript{(39)} を閉め、\VUgras{窓の帳} \textsuperscript{(41)} を巻き上げ、
小さな包を\VUgras{荷棚} \textsuperscript{(43)} に載せてから、\VUgras{座席}
\textsuperscript{(42)} に腰を下ろした。\VUgras{灯係} \textsuperscript{(44)} が灯を
点けるのを見て、車の中で一夜を過ごすのだと思い出し、\VUgras{枕}
\textsuperscript{(45)} を貸す係を呼んで一つもらった。

%%K t12-11-1 p
11. --- 車掌が切符を調べに来た。時刻は過ぎているのになぜ発たぬのかと尋ねた。
\VUgras{車長} \textsuperscript{(46)} が自転車を\VUgras{荷の車輌} \textsuperscript{(47)}
に積ませているところだと答えた。そのうえ\VUgras{足温器} \textsuperscript{(48)} も
まだみな替え終わっていないという。

%%K t12-12-1 p
12. --- だがじきに発つはずであった。\VUgras{火夫} \textsuperscript{(50)} が自分の
\VUgras{炭水車} \textsuperscript{(49)} の上で石炭を取り、\VUgras{機関車}
\textsuperscript{(54)} の火を強めていたし、\VUgras{機関手} \textsuperscript{(51)} は
\VUgras{助役} \textsuperscript{(52)} の合図を待つばかりであった。助役は
\VUgras{ホーム} \textsuperscript{(53)} に立っていた。

%%K t12-13-1 p
13. --- 機関車の\VUgras{汽罐} \textsuperscript{(56)} が低く唸り、\VUgras{煙突}
\textsuperscript{(55)} が\VUgras{蒸汽} \textsuperscript{(60)} の混じった煙を吐いた。
鋭い\VUgras{汽笛} \textsuperscript{(59)} が鳴り、\VUgras{ピストン}
\textsuperscript{(58)} が動き出して、この重い機械の\VUgras{車輪} \textsuperscript{(57)}
を回した。

%%K t12-tit-3 sub
\VUsaut{3.0mm}
\VUpk{10.2pt}{40}{発車！}
\VUsaut{3.0mm}

%%K t12-14-1 p
14. --- 汽車は\VUgras{軌条} \textsuperscript{(64)} の上を走り出し、速さを増した。
\VUgras{ホーム} \textsuperscript{(65)} は尽き、\VUgras{便所} \textsuperscript{(66)} の
そばを過ぎ、別の\VUgras{線} \textsuperscript{(63)} に置かれた車輌の列のそばも過ぎた。
\VUgras{寝台車} \textsuperscript{(67)}、\VUgras{食堂車} \textsuperscript{(68)}、
\VUgras{郵便車} \textsuperscript{(69)} である。

%%K t12-noto-1 noto
\VUnotes{143.2mm}{%
\textit{(*) 注。} --- \textit{この旅は言うまでもなく戦前の国境の有様に従って
書いたものである。}%
}

%%K t12-15-1 p
15. --- 次に\VUgras{貨物の停車場} \textsuperscript{(71)} が目の前を過ぎた。覆屋の
下では役人たちが緩行で託された\VUgras{荷} \textsuperscript{(72)} を積んでいた。
\VUgras{入換の人} \textsuperscript{(76)} が\VUgras{水槽} \textsuperscript{(73)} の
そばで\VUgras{家畜車} \textsuperscript{(75)} と\VUgras{無蓋車} \textsuperscript{(79)}
を動かし、一列の\VUgras{貨物列車} \textsuperscript{(80)} を組んでいた。

%%K t12-16-1 p
16. --- 我々は最後の\VUgras{転轍} \textsuperscript{(78)} を過ぎた。そばに転轍手が
立っていた。\VUgras{信号機} \textsuperscript{(86)} も後にし、停車場は遠くに消えた。

%%K t12-17-1 p
17. --- まもなく我々はあの\VUgras{鉄の橋} \textsuperscript{(74)} に乗った。橋は
\VUgras{川} \textsuperscript{(81)} を跨いでいる。橋を渡ると川に沿って走った。川では
力強い\VUgras{曳船} \textsuperscript{(82)} が重い\VUgras{荷足船} \textsuperscript{(83)}
を曳いていた。\VUgras{客船} \textsuperscript{(84)} が\VUgras{浮桟橋}
\textsuperscript{(85)} に着こうとするところも見えた。

%%K t12-18-1 p
18. --- 幾つもの停車場を通り過ぎたのち、幾つかの\VUgras{隧道} \textsuperscript{(87)}
を抜け、\VUgras{踏切番}の数えきれぬ\VUgras{小屋} \textsuperscript{(88)} を後に
していった。車の揺れに揺られて、私は深く眠りこんだ。

%%K t12-tit-4 sub
\VUsaut{3.0mm}
\VUpk{10.2pt}{40}{ドイチュ・アヴリクール停車場。}
\VUsaut{3.0mm}

%%K t12-19-1 p
19. --- 私は突然、役人の叫び声に起こされた。「ドイチュ・アヴリクール。
\textsuperscript{(*)} 皆さんお降り。」税関の検査と、国境のあらゆる面倒な手続で
ある。私は停車場の\VUgras{時計} \textsuperscript{(89)} に自分の時計を合わせねば
ならなかった。その時計は五十五分進んでいた。眠い目で私は\VUgras{長針}
\textsuperscript{(90)} と\VUgras{短針} \textsuperscript{(91)} をやっと見分けるほどで
あった。\VUgras{文字盤} \textsuperscript{(92)} そのものが朝の霧でかすんでいた。

%%K t12-20-1 p
20. --- 税関の者が検めているあいだ、私は欠伸をしながら停車場の壁を飾る
\VUgras{はり紙} \textsuperscript{(93)} を読んでいた。暇つぶしに朝食をとり、ドイツの
\VUgras{鉄道網} \textsuperscript{(94)} の図でストラスブールまでどれほどあるかを見、
旅の終りが近いという望みに元気づけられて、また車輌に戻った。
\VUsaut{6.0mm}
\VUfilet{22mm}
